\section{Introduction}
\label{sec:intro}

The Riemann hypothesis asserts that the non-trivial zeros of the function
$\zeta(s) = \sum_{n \ge 1} n^{-s}$ have real part $1/2$. Beyond its
arithmetic formulation, this conjecture is also a question about the
spectral nature of a certain operator, whose existence has been
hypothesised since Hilbert and P\'olya. \textcite{BerryKeating1999} in 1999,
then \textcite{Connes1996} in 1996, made the content of this spectral
hypothesis precise: it is the dilation operator
$H = (xp + px)/2$ quantised on the half-line, whose microcanonical
semi-classical quantisation reproduces the Riemann--von Mangoldt density
$N(\gamma) \sim (\gamma/2\pi) \log(\gamma / 2\pi e)$
of the zeros.

The programme presented here does not seek a proof of the Riemann
hypothesis. It seeks to reconstruct the Hilbert--P\'olya--Berry--Keating
(HP-BK) operator framework canonically from a single internal structure,
that of the Theory of Persistence (PT), which derives the Standard Model
from a unique arithmetic parameter $\sper = 1/2$
(theorem~$\Tone$, \cite{PT_M1}). The work revisits the RH question from
three complementary angles.

\paragraph{First angle: analytic.}
We exhibit a factorisation
$\zeta(s) = \Rres(s)\,\zetaplus(s)\,\zetaminus(s)$, where $\zeta_\pm$ are
defined by the amplitudes $a_p, b_p$ coming from the $\qplus/\qminus$
bifurcation of PT. The residual factor $\Rres(s)$ admits an exact Fredholm
representation $\det(I - \DRop(s))^{-1} = \Rres(s)$ with $\DRop$ diagonal
and eigen-amplitudes $\kappap(s)$ in closed form (theorem~\ref{thm:fredholm}
in \S\ref{sec:fredholm}). This representation is proved in $\Reop(s) > 1$
and verified numerically to $10^{-12}$.

\paragraph{Second angle: operator-theoretic.}
We identify the canonical operator
$\HPTBK = (u\,p_u + p_u\,u)/2$ on the Mellin coordinate $u = \log p$
with its dynamical cut-off
$\umax(\gamma) = \gamma/\sqrt{2\pi}$. This cut-off is not postulated; it
is forced by the symplectic Planck cell $u_\star\,p_\star = 2\pi$
of the canonical phase space (\S\ref{sec:cellule_planck}). We prove that
the microcanonical semi-classical quantisation of this operator, in
symmetric Weyl ordering with double barrier, reproduces the exact
functional form of the Riemann--von Mangoldt density, up to a constant
residue $1/8$, verified to $10^{-15}$ on four decades
$\gamma \in [14, 10^{3}]$ (\S\ref{sec:densite_semi_classique}).

\paragraph{Third angle: spectral and distributional.}
The construction of a regularised trace
$\Treg(s) = \Tr_{\mathrm{reg}}\!\big[(s - i\HPTBK)^{-1}\,\DRop(s)\big]$
links the canonical operator to the residual factor via the explicit trace
formula
\[
-\dfrac{d \log \Rres}{ds} \;=\; \sum_{p} (\log p)\,\Big[\,\dfrac{1}{p^s - 1} \,+\, \Ap\,p^{-s}\,\Big],
\]
proved as an exact analytic identity in $\Reop(s) > 1$
(theorem~\ref{thm:trace_formula}) and distributionally continued to the
critical line (\S\ref{sec:trace_reg}). The Hadamard regularisation of this
trace produces poles at the ordinates $\gamma_n$ of the zeros of
$\Rres$; we capture $17$ zeros out of $17$ tested on $t \in [10, 70]$ at
$|\Delta| < 1.5$, and $74$ out of $79$ on $t \in [10, 200]$ ($94\,\%$),
with saturation beyond $P_{\max} = 3\,000$.

\paragraph{Intrinsic enrichment.}
The programme distinguishes itself from Berry--Keating--Connes by the
intrinsic enrichment it provides. The apparently geometric numbers of the
classical construction --- the $2\pi$ of the Planck cell, the $1/8$ of
the Maslov phase of the corners (semi-classical WKB correction of $-\pi/8$
per turning point at the corners of the microcanonical area domain,
cf.~\cite[\S 3]{BerryKeating1999}) --- receive a derivation proper to
PT. The $2\pi$ is the canonical Liouville measure on $(u, p_u)$,
that is the commutator $[u, p_u] = i$. The $1/8$ admits a triple
cohomological reading:
\[
\dfrac{1}{8}
\;=\; \dfrac{\sper^{2}}{2} \;=\; \lambdaH
\quad\text{(Higgs self-coupling, \cite{PT_M5})},
\]
\[
\dfrac{1}{8}
\;=\; \dfrac{c_1}{N_{\mathrm{corners}}}
\quad\text{(first Chern class of the $\qplus/\qminus$ K\"ahler--Fisher spinor bundle)},
\]
\[
\dfrac{1}{8} \;=\; \dfrac{\sper}{4}
\quad\text{(Berry phase per corner)}.
\]
The consistency of the three readings reduces to the equation
$\sper^{2}/2 = \sper/4$, whose only solutions are $\sper = 0$ and
$\sper = 1/2$; axiom $\Tone$ excludes $\sper = 0$, which overdetermines
the PT parameter modulo $\Tone$ assumed (and not as an independent
derivation from $\Tone$).

\paragraph{Honest falsification: Dirichlet.}
A prediction stemming from this framework, the formula
$\Delta_N(L \bmod q) = 1/(8q)$ for primitive Dirichlet $L$-functions, has
been numerically tested on \texttt{LMFDB} zeros and falsified
(\S\ref{sec:auto_correction}). This falsification is explicitly
documented. It leads to reformulating the status of the $1/8$ established
in \S\ref{sec:densite_semi_classique}: it is an asymptotic identity
between smooth forms, not a measurable residue in the count of zeros. The
PT mechanism described here is specific to $\zeta$, and not to the
Dirichlet family via a naive factorisation by the Haar measure
modulo~$q$.

\paragraph{Programme position.}
The distinction maintained throughout the text is the following. We seek
to understand the structure of the zeros, to provide a canonical framework
without ad hoc parameter, to derive the known numerical coincidences
($2\pi$, $7/8$, critical line) from a single internal structure. We do
not seek to prove that all zeros lie on the critical line: this ultimate
lock (equivalent to RH in the strict sense) remains open and does not
seem easier than the classical version.

